% Source: https://artofproblemsolving.com/community/c3h1590040_two_troll_mc_geo_problems
% Stable identity: aops-topic-1590040-q02
\textbf{The removed cone.} The removed top cone is similar to the original one (the cut is parallel to the base), with some ratio $k$. Its volume is $k^3$ times the original volume, and it must be $1-\dfrac{1}{3}=\dfrac{2}{3}$ of it, so
\[
k^3=\dfrac{2}{3},\qquad k=\left(\dfrac{2}{3}\right)^{\dfrac{1}{3}}.
\]
So the top cone has radius $5k$ and slant height $10k$, and the frustum has radii $5$ and $5k$ and slant height $10-10k$:

\begin{asy}
usepackage("amsmath");
size(4.8cm);
defaultpen(fontsize(10pt));
real R=5,H=5*sqrt(3),k=(2/3)^(1/3),ey=0.2;
pair V=(0,H);
real hc=H*(1-k), rc=R*k;
path el(pair c,real r){return shift(c)*scale(r,ey*r)*unitcircle;}
fill((-R,0)--(-rc,hc)--(rc,hc)--(R,0)--cycle,palegray);fill(subpath(el((0,0),R),2,4)--cycle,palegray);
draw(subpath(el((0,0),R),0,2),dashed);
fill(el((0,hc),rc),gray(0.85));draw(subpath(el((0,0),R),2,4),linewidth(0.9));
draw(el((0,hc),rc),linewidth(0.9));
draw((-R,0)--(-rc,hc),linewidth(0.9));draw((R,0)--(rc,hc),linewidth(0.9));
draw((-rc,hc)--V--(rc,hc),gray+dashed);
dot(V,gray);label("$V$",V,N,gray);
draw((0,0)--(R,0),blue);label("$5$",(R/2,0),S,blue);dot((0,0));
pair n=unit(rotate(90)*(V-(-R,0)));
draw((-R,0)+0.8*n--V+0.8*n,red,Arrows(TeXHead));
label("$10$",(-R,0)+0.55*(V-(-R,0))+0.8*n,W,red);
label("\mbox{}",(min(currentpicture,user=true)+max(currentpicture,user=true))/2,black); // keeps the text after the figure black when LaTeX typesets the labels inline
\end{asy}

(The removed top cone, with apex $V$, is dashed.)

\textbf{The lateral area.} The lateral area of a cone with radius $r$ and slant height $\ell$ is $\pi r\ell$. The frustum's lateral area is the difference of two such areas:
\[
\pi\cdot5\cdot10-\pi\cdot5k\cdot10k=50\pi\left(1-k^2\right).
\]

\textbf{The total area.} Adding the two circular bases, of areas $25\pi$ and $25\pi k^2$,
\[
50\pi\left(1-k^2\right)+25\pi+25\pi k^2=75\pi-25\pi k^2=25\pi\left(3-\left(\dfrac{2}{3}\right)^{\dfrac{2}{3}}\right).
\]
Numerically $\left(\dfrac{2}{3}\right)^{\dfrac{2}{3}}\approx0.76314$, so the area is about $25\pi\cdot2.23686\approx175.68$, which rounds to $\boxed{176}$ square meters.
